\section{Conclusion}

In this work, we revisit hallucination in large language models from a behavioral perspective. 
Rather than treating hallucination as a purely generative error, we formulate it as a decision-making problem under uncertainty, involving when to answer, how to answer, and how confidently to respond.

We propose a behavior-oriented reinforcement learning framework that integrates behavior alignment to distinguish between answering and abstaining, uncertainty modeling via entropy-based overconfidence control, response quality shaping to ensure usefulness and completeness, and length regularization to prevent degenerate generation strategies.

Through extensive experiments, we demonstrate that the proposed approach reduces hallucination by over 65\%, consistently outperforming strong baselines such as GPT-4o. 
Furthermore, the learned behavior generalizes effectively to unseen domains, significantly improves response quality across multiple dimensions, and preserves general knowledge capability without catastrophic forgetting.

These results suggest that hallucination is not merely a limitation of knowledge, but a consequence of suboptimal decision-making under uncertainty. 
By explicitly modeling behavior, confidence, and response quality, the proposed framework enables language models to produce outputs that are both reliable and useful.